\begin{longtable}{p{0.20\textwidth}p{0.33\textwidth}p{0.38\textwidth}}
\toprule
Artifact & What was checked & What the check does not establish\\
\midrule
\endhead
Reflection kernel & Pinned Lean accepted the all-offset theorem, sparse-support adapter, sum lifting, wrappers and boundary controls. & The compatibility-preserving production patch was source-reviewed; the complete imported project module and its clients were not rebuilt.\\
Prime-support kernel & Pinned Lean accepted fixed-residue unit transfer and the guarded polynomial congruence, including exponent zero and a composite thickened factor. & It does not evaluate the literal current kernel or discharge the concrete fixed unit certificate.\\
Subset-translation kernel & Pinned Lean accepted the additive-group theorem, the modular unit instance, and the existence endpoint. & It does not replace the public exact interval definitions or establish safe physical mutation.\\
Prime-power replacement & Complete original/replacement source, existing checks, added arbitrary/zero-exponent controls, and axiom comparisons all passed. & This is a focused module replay, not a whole-repository build or a downstream auxiliary-order client rebuild.\\
\bottomrule
\end{longtable}

\noindent The prime-power patch preserves all six original public statement
texts and axiom sets, and adds \code{lift_divisibility}. Including the new
helper, the full file changes from 80 to 77 physical lines and from 3,415
to 3,425 bytes. The reflection proposal preserves the old source theorem
as a wrapper and changes its full module from 350 to 351 lines; its shorter
core proof is not advertised as a net source deletion. No net deletion is
claimed for the two larger architectural prototypes. The logged compiler
runs contain no \code{sorryAx}; standard Lean/Mathlib axioms are recorded
in the receipts. Linter diagnostics in concrete reflection examples do not
affect kernel acceptance.
